\section{Fundamentals and Related Work}
\label{sec:fundamentals}

This section defines urban accessibility, describes how it is computed over a travel network, and reviews the literature on accessibility under environmental disruption and on spatial inequality.

\subsection{Urban Accessibility}

In urban and transport studies, urban accessibility describes how easily opportunities — services, jobs, facilities, and activities — can be reached from an origin. Unlike mobility, which is observed movement, accessibility is a structural property of the urban system: it depends on where opportunities are located, on how the network connects origins and destinations, and on which impedances (costs) are imposed on travel~\cite{Handy1997Measuring}. Accessibility is therefore usually treated as a measure of access potential rather than a direct measure of realized trips.

Formally, the territory is represented by a set of origins $O$ — territorial cells, tract centroids, blocks, or residence points — and a set of destinations $D$ (points of interest, POIs). A cost function $c(o,d)$ assigns to each pair $(o,d)$ the minimum travel cost between $o \in O$ and $d \in D$~\cite{Miller1999Measuring}. In network-based applications this cost is typically obtained by shortest-path algorithms over a travel graph, using walking or cycling time, and it provides the basis for indicators that are comparable across areas and scenarios.

Although often expressed as time or distance, $c(o,d)$ is better understood as a composite impedance that also captures comfort, safety, physical effort, and route reliability~\cite{Puello2016Integration}. This matters for short-distance active mobility~\cite{Cook2022More}, where the experience of the route and local conditions — slope, crossings, sidewalk quality, exposure to sun, rain, or wind, and the presence of flooding — can alter the viability of a trip even when the network geometry is unchanged. Climate disruptions can thus enter the problem as changes to the generalized travel cost, rather than only as binary segment closures.

Within location-based accessibility measures, reviews such as \cite{Geurs2004} group indicators by their assumptions: distance or minimum-cost measures, cumulative (contour) measures, and potential (gravity) measures. These families differ in how they treat destination choice, substitutability among opportunities, and distance or time decay, yet they share the premise that accessibility is derived from an origin-destination cost matrix.

The simplest measures are distance metrics, in which the accessibility of an origin is the minimum cost to the nearest opportunity — for example, the time to the nearest shop. Such an indicator is intuitive and adequate when users are assumed to satisfy a need at the easiest destination to reach. It ignores the existence of multiple alternatives, does not represent competition among offers, and can be unstable when small network changes alter which POI is nearest.

When the availability of multiple opportunities matters, cumulative (contour) measures count how many opportunities are reachable within a cost threshold, such as the number of POIs within 15 minutes. A useful variant that avoids fixing a threshold while retaining an intuitive reading is to compute accessibility from the mean cost of reaching the $N$ nearest opportunities in each service category. The metric then represents the average ease of reaching a minimum set of relevant alternatives, which matches short-distance travel where residents effectively consider only a few options.

Despite their clarity, cumulative and $N$-nearest metrics carry operational choices that affect the result: the threshold value where one exists, the value of $N$, sensitivity to ties, and whether opportunities are weighted equally. These measures also tend to treat cost as an abrupt cut — inside or outside the reachable set — or as a simple average, which can under-represent gradual changes in destination attractiveness as cost increases. Under climate disruption, where cost may change continuously and unevenly across space, formulations that progressively penalize distant destinations without a hard cutoff become preferable.

In potential (gravity) metrics, accessibility is the sum of opportunities weighted by a cost-decay function, capturing the idea that distant destinations contribute less than nearby ones. The classical formulation of \cite{Hansen1959} expresses the accessibility of an origin $o$ as
\begin{equation}
A(o) = \sum_{d \in D} W_d \, f\!\big(c(o,d)\big),
\end{equation}
where $W_d$ is the mass or attractiveness of the opportunity at destination $d$ — capacity, size, or a unit weight — and $f\!\big(c(o,d)\big)$ is a decreasing function that reduces the contribution as cost increases. In practice the destination set is restricted to a relevant neighborhood, such as the $N$ nearest destinations or those within a cost window, which preserves the decay principle while preventing very distant opportunities from having undue influence and keeping the indicator computationally tractable.

The decay function $f(c)$ may be exponential, power, or logistic, and its parameters control how steeply accessibility falls as cost increases. Conceptually this approximates different choice behaviors and travel tolerances; operationally, the parameters can be set from the literature, from observed data, or through sensitivity analysis. For short-distance active mobility, relatively steep decay is consistent with the rapid drop in willingness to walk or cycle as cost rises, especially when cost incorporates environmental discomfort.

In the broader conceptual frame, accessibility results from the interaction of land use (where opportunities are), transport (network and impedances), temporality (variation by hour, season, or event), and individual characteristics (constraints and preferences). This work focuses on the transport and impedance component, comparing a typical scenario, in which costs derive from network topology and usual weights, against a scenario conditioned by climate disruption, in which edge weights are penalized by environmental variables. Holding opportunities and the spatial representation of origins constant, the difference between indicators across the two scenarios expresses the sensitivity and robustness of accessibility under climate disruption.

\subsubsection{Short-Distance Accessibility Indicators}
\label{subsubsec:indicadores}

To compute urban accessibility to everyday services, \cite{Bruno2024} use a set of territorial and inequality indicators over a regular spatial grid as the minimum unit of analysis.

\begin{description}
  \item[PTh (Proximity Time per territorial unit).] Mean time required to reach the $k$ nearest points of interest of a service category \ric{$i$}? from each territorial cell:
\begin{equation}
    PTh_i = \frac{1}{k} \sum_{d \in D_k(i)} c(i,d),
\end{equation}
\noindent where $D_k(i)$ is the set of the $k$ nearest POIs and $c(i,d)$ is the minimum travel time over the network. The parameter $k$ is adjustable: $k=1$ for emergency health care, $k=2$ or $k=3$ for education and retail. Lower PTh values indicate greater accessibility. Computed per cell, PTh forms a spatial layer that reveals territorial gradients and patterns of opportunity concentration.

  \item[G index (accessibility inequality).] Gini coefficient measuring how accessibility is distributed across territories or socioeconomic groups. It ranges from 0 (perfect equality) to 1 (maximum inequality). Higher values indicate greater concentration of poor accessibility in specific segments of the population. The index is particularly useful for comparing how the accessibility distribution shifts between scenarios, revealing whether environmental disruptions amplify or reduce inequalities.

  \item[F15 (share of the population in 15-minute zones).] Proportion of the resident population in cells where PTh $\leq$ 15 minutes:
\begin{equation}
    F15 = \frac{\sum_i p_i \cdot \mathbb{1}(PTh_i \leq 15)}{\sum_i p_i} \times 100\%,
\end{equation}
\noindent where $\mathbb{1}(\cdot)$ is an indicator function and $p_i$ is the resident population in territorial unit $i$. Unlike PTh, F15 is binary — inside or outside the threshold — and complements it in assessing population coverage under the 15-minute city model.
\end{description}

\subsection{Mobility Graphs and Networks}

The minimum cost $c(o,d)$ is generally obtained from a travel network represented as a weighted graph $G=(V,E)$. Vertices $V$ correspond to decision points in the street network, such as intersections and topological nodes, while edges $E$ represent traversable segments: streets, passages, and connections. Each edge $e \in E$ carries a weight $w(e)$ that encodes the effort required to cross it.

The network is often treated as undirected, assuming that travel along a segment costs the same in both directions. Some situations, however, make impedance direction-dependent, which calls for a directed graph that allows distinct costs per orientation of the same connection — because of slope, stairs, controlled crossings, or circulation restrictions. This distinction changes the set of feasible paths considered when computing $c(o,d)$.

For walking, a simple and widely used model defines $w(e)$ as travel time, estimated from the segment's geometric length and an assumed walking speed. In more general formulations the weight is a generalized cost that incorporates penalties related to environmental quality. The cost between an origin $o$ and a destination $d$ is then the cost of the shortest path in the network, that is, the minimum sum of weights along a feasible path:
\begin{equation}
    c(o,d) = \min_{p} \sum_{e \in p} w(e),
\end{equation}
with $p \in \mathcal{P}(o,d)$, the set of feasible paths between $o$ and $d$. Algorithmically, this value is obtained by classical shortest-path methods such as Dijkstra's algorithm~\cite{cormen01introduction} or heuristic variants such as $A^*$, depending on graph properties and performance requirements.

Changes in travel conditions are represented in the network either as continuous impedance changes or as circulation restrictions that make certain segments unavailable. Both modify the feasible paths between origins and destinations, and therefore the estimated minimum costs. A pair $(o,d)$ may even become unreachable.

Evaluating origins and destinations on the network also requires defining how these points connect to the graph. In urban applications, POIs and origins given by zone or territorial-unit centroids are typically snapped to the nearest node or segment, or connected by artificial edges that represent local access to the network. Though operational, this step influences the estimated minimum costs and must be consistent across scenarios to preserve comparability.

\subsection{Related Work}

\subsubsection{Network- and POI-Based Accessibility}

This work belongs to the line of research that evaluates urban accessibility from travel networks and the spatial distribution of opportunities. In this direction, \cite{Bruno2024} propose a methodology to estimate access times to everyday points of interest over a hexagonal spatial grid, producing local indicators per territorial unit and city-level aggregate metrics, including mean access time to the nearest points of interest and population-weighted summaries.

Gravity metrics are common in urban accessibility analysis, often using the number of job opportunities in a territorial unit as that area's attractiveness. \cite{PRITCHARD2019386} developed a potential accessibility model based on the gravity form of \cite{Hansen1959}, using a negative exponential decay function weighted by opportunities. This approach captures travel behavior by giving greater weight to nearby opportunities, and departs from cumulative measures that treat all trips within a time interval as equally accessible. The methodology adopted here follows a similar approach but extends the notion of generalized cost to incorporate climate disruptions as continuous impedance factors.

\subsubsection{Accessibility under Environmental Disruption}

Studies that assess the impacts of disruptions commonly model network degradation by removing impassable segments and quantifying the resulting fragmentation and increase in travel costs. \cite{Loreti2022} analyze the road network between urban centers, removing segments compromised during floods and evaluating how this affects accessibility between localities, which allows identifying situations of isolation. \cite{Abenayake2022An} propose a methodology to assess transport system failures by comparing a non-flooded reference condition against actual flood events. The method overlays the road network on flooded areas and assumes that affected segments become inaccessible, recomputing centrality metrics (closeness and betweenness) to characterize topological and functional changes in network performance. At the intra-urban scale, \cite{Bono2011A} combine graph concepts and spatial analysis in \ac{SIG} to model post-disaster road interruptions (earthquakes), identify fragmentation into connected components, and estimate accessibility reduction by comparing travel costs obtained for the unmodified and the damaged network.

In contrast to these binary-removal approaches, recent studies have shown that climate disruptions often increase travel costs continuously and unevenly across space without necessarily removing network segments. An important methodological contribution is the modeling of spatial variability in travel costs. \cite{PRITCHARD2019386} incorporate topography (slope), infrastructure, waiting times, and congestion to estimate realistic cyclist and pedestrian speeds across segments. Similarly, this work models how climate disruptions increase travel costs through penalty functions that reflect the degradation of the pedestrian experience, rather than simply removing network segments.

\subsubsection{Equity and Inequality in Accessibility}

Urban accessibility is traditionally evaluated by metrics that privilege aggregate network efficiency and disregard socio-spatial inequalities. \cite{FolGallez2014} show that this approach fails to capture barriers faced by vulnerable populations. Recent studies therefore adopt coverage measures within 15-minute zones combined with inequality indicators such as the Gini index \cite{Bruno2024}. This work follows that direction by examining intra-urban clusters with similar access times and socioeconomic conditions, investigating how climate disruptions amplify pre-existing inequalities.

Recent studies have begun to integrate environmental analysis with equity within 15-minute zones. \cite{smartcities8020053} examine disparities in air pollution exposure within 15-minute zones in New York, revealing how low-income and Black communities experience significantly higher pollutant levels in their walking ranges, which demonstrates the importance of integrating network analysis with socioeconomic equity.

Along the same lines, studies have begun to investigate how acute climate disruptions interact with social vulnerability within 15-minute zones. \cite{Shartova2024} investigated thermal vulnerability in Moscow during the 2021 heat waves, revealing that areas with high thermal stress do not necessarily coincide with high social vulnerability, and that new residential zones with little green infrastructure are the most vulnerable. None of these works, however, examines how climate disruptions degrade accessibility to essential services by comparing scenarios, nor incorporates continuous cost functions that reflect how pedestrians experience environmental degradation — gaps this work addresses.

To measure inequalities, \cite{PRITCHARD2019386} use the Gini coefficient and apply clustering techniques to identify regions with different accessibility levels. A critical finding is that improvements in bicycle accessibility concentrate in middle- and high-income areas, leaving peripheral and poor areas with minimal gains. This suggests that environmental disruptions such as extreme climate events likely amplify pre-existing inequalities, a premise this work investigates.

\subsubsection{Comparative Synthesis of the Literature}

\hide{\ric{Precisa citar e explicar a Tabela~\ref{tab:trabalhos_relacionados} no texto. Por exemplo, a Tabela~\ref{tab:trabalhos_relacionados} resume alguns dos trabalhos relacionados segundo o tipo de interdição da malha viária causada pela condição ambiental. Neste caso, pode haver a simples interdição ou uma penalização de caminhos... Além disso, os trabalhos podem considerar ou não a condição econômica da população afetada.}}

Table~\ref{tab:trabalhos_relacionados} summarizes the reviewed works by the type of environmental effect modeled on the road network, their treatment of typical and conditioned scenarios, and whether they include socioeconomic analysis. Much of the literature represents environmental disruptions as binary interdiction of network segments, producing abrupt cuts in connectivity and accessibility. Other studies model continuous penalties on travel cost, capturing gradual variation without complete disconnection of the graph. Explicit association between accessibility differences, environmental conditions, and socioeconomic variables remains infrequent. This work advances by combining continuous cost penalization, comparative scenario analysis, and integration with socioeconomic indicators, extending the capacity to identify territorial vulnerabilities under climate disruption.

\begin{table}[H]
\centering
\footnotesize
\caption{Comparison of related work on urban accessibility, disruption, and equity}
\label{tab:trabalhos_relacionados}
\begin{tabularx}{\linewidth}{|>{\centering\arraybackslash}p{3.7cm}|>{\centering\arraybackslash}X|>{\centering\arraybackslash}p{3cm}|>{\centering\arraybackslash}X|>{\centering\arraybackslash}X|}
\hline
\textbf{Reference} &
\textbf{Network/graph} &
\textbf{Environmental effect on network} &
\textbf{Typical vs. conditioned} &
\textbf{Socioecon. analysis} \\
\hline
\cite{Bruno2024}           & Yes          & N/A                   & No           & No \\\hline
\cite{Loreti2022}          & Yes          & Interdiction          & Yes          & No \\\hline
\cite{Abenayake2022An}     & Yes          & Interdiction          & Yes          & No \\\hline
\cite{Shartova2024}        & No           & N/A                   & No           & Yes \\\hline
\cite{smartcities8020053}  & Yes          & Exposure              & No           & Yes \\\hline
\cite{Bono2011A}           & Yes          & Interdiction          & Yes          & No \\\hline
\cite{PRITCHARD2019386}    & Yes          & Penalization          & No           & Yes \\\hline
\textbf{This work}         & \textbf{Yes} & \textbf{Penalization} & \textbf{Yes} & \textbf{Yes} \\\hline
\end{tabularx}

\vspace{0.4em}
\begin{minipage}{0.98\linewidth}
\footnotesize
\textit{Notes:}
% (1) \textbf{Interdição}: a condição ambiental remove/fecha segmentos, alterando a conectividade do grafo.
% (2) \textbf{Penalização}: a condição ambiental repondera continuamente o custo das arestas (ex.: aumento de tempo/esforço), afetando rotas e acessibilidade sem necessariamente desconectar a rede.
(1) \textbf{Exposure}: the environmental condition is a continuous attribute attached to segments and aggregated as an exposure metric within a reachable area (e.g., 15 min), without altering the cost or travel time used for routing.
(2) \textbf{N/A}: the study does not incorporate an environmental condition into network modeling.
\end{minipage}
\end{table}
